La gestión de archivos dentro del sistema se divide en dos
fases principales: la lectura de los documentos de origen
(Entrada) y la generación del reporte final (Salida).

Para la etapa de entrada, los archivos fuente se ubican
en un directorio interno del proyecto \newline
(\texttt{src/main/resources/pdf/}).
Desde esta ruta, los documentos son leídos y procesados
de acuerdo con las lógicas descritas previamente en las
secciones \ref{sec:hilos} Manejo de hilos y 
\ref{sec:algoritmos} Algoritmos de prioridad. 

Las rutas de lectura y salida predeterminada están
declaradas como constantes estáticas dentro de la clase
\texttt{Generador.java}, lo que facilita el mantenimiento
del código:

\vspace{-0.2cm}
\begin{figure}[H]
    \centering
    \caption{Declaración de las rutas de entrada y salida predeterminada}
    \label{code:path_entrada_archivos}
    \vspace{-0.2cm}
    \begin{lstlisting}[language=Java]
public class Generador {
    private static final Path CARPETA_PDF =
            Paths.get("src/main/resources/pdf/");
            
    private static final Path ARCHIVO_SALIDA =
            Paths.get("horarios.txt");
    // ...
}
    \end{lstlisting}
\end{figure}

Posteriormente, estas constantes son utilizadas por el método
\texttt{cargarAlumnosDesdePDF()} para iniciar el
procesamiento de cada documento.

En cuanto a la etapa de salida, una vez que la información
ha sido extraída exitosamente, el sistema procede a escribir
los resultados. Aunque existe un nombre de archivo por defecto
(\texttt{horarios.txt}), el sistema no guarda el documento
de forma arbitraria. Para ofrecer una mejor experiencia,
el programa delega la selección de la ruta de guardado
a la clase \texttt{SelectorDestino.java}. Esta clase hace
uso de la herramienta gráfica \texttt{JFileChooser} para
abrir un cuadro de diálogo nativo del sistema operativo,
permitiendo al usuario elegir manualmente el directorio
de destino final, tal y como se detalló en la sección
\ref{sec:SelectorDestino} Clase SelectorDestino.